\documentclass[12pt]{book}
\usepackage[utf8]{inputenc}
\usepackage[T1]{fontenc}
\usepackage{amsmath,amssymb,amsfonts}
\usepackage{amsthm}

\begin{document}

\setcounter{page}{326}
which is of finite type over \(V\).
Then \(U\) admits a closed immersion \(i\) into an affine space \(\mathbb{A}_{V}^{n}\) for suitable \(n\).
Let \(P=\mathbb{A}_{V}^{n}\), let \(p: P \to Y\) be the natural projection, let \(f^{a}=i^{y} p^{z}\), and \(\gamma_{i,p}=id\).
This gives a chart on \(U\) which is a neighborhood of the given point \(x\).

If \(f^{a}\) is a chart on an open set \(U\), and if \(U' \subseteq U\) is a smaller open set, we define the notion of a restriction of the chart to \(U'\), as follows: let \(P' \subseteq P\) be an open set whose intersection with \(i(U)\) is \(i(U')\).
Then we take \(f^{a}|_{U'}\), \(i|_{U'}\), \(P'\), \(\gamma_{i,p}|_{U'}\), as the restriction of the chart.
Note that a restriction of permissible isomorphisms is permissible.

For each pair of indices \(\mu,\nu\) choose restrictions of the charts on \(U_{\nu}\) and \(U_{\mu}\) to \(U_{\mu\nu}=U_{\nu}\cap U_{\mu}\), and let \(\alpha_{\mu\nu}\) be the unique permissible isomorphism between them.
One sees immediately that the isomorphism of functors

\setcounter{page}{327}
\[\alpha_{\mu\nu}: f^{\nu} \to f^{\mu}\]
thus defined on \(U_{\mu\nu}\) is independent of the restrictions of the charts chosen.
Furthermore it follows from the lemmas that on a triple intersection \(U_{\mu\nu\lambda}\) these isomorphisms are compatible:
\[\alpha_{\mu\nu}\alpha_{\nu\lambda}\alpha_{\lambda\mu}=id .\]

Thus, since we are dealing with functors on residual complexes, we can use Lemma 1.3, and glue the functors \(f^{\nu}\) via the isomorphisms \(\alpha_{\mu\nu}\) to obtain a functor
\[f^{\Delta}: Res(Y) \to Res(X)\]
(together with isomorphisms \(\beta_{\nu}:f^{\Delta}|_{U_{\nu}} \to f^{\nu}\) compatible with \(\alpha_{\mu\nu}\)) which is the one we want.

b) Construction of the isomorphisms \(c_{f,g}\).
Let \(f: X \to Y\) and \(g: Y \to Z\) be morphisms.
It will be sufficient to construct \(c_{f,g}\) locally, chart by chart, and show that it is compatible with the permissible isomorphisms of change of chart.
For then we can glue.
Thus we may assume that \(f\) and \(g\) are embeddable in smooth morphisms, and we may even assume that \(X\) is embeddable in an affine space over \(Y\), since that

\setcounter{page}{328}
so that we have a commutative diagram as shown, with \(i,j,k\) closed immersions, \(p,q,r\) smooth, and the upper right square Cartesian.
Define \(c_{f,g}\) as follows:
\[c_{f,g}:(g f)^{a} \stackrel{\gamma_{ji,rq}}{\to}(j i)^{y}(r q)^{z} \stackrel{I,\Pi}{\to} i^{y} j^{y} q^{z} r^{z} \leftarrow \Pi_{i} y_{p} z_{k} y_{r} z \stackrel{\gamma_{i,p}\gamma_{k,r}}{\leftarrow} f^{b} g^{c} .\]

Now we must show that \(c_{f,g}\) is compatible with the permissible isomorphisms of change of chart.
It is sufficient to vary one at a time.
Furthermore, if \(k': Y \to Q\) is another chart for \(g\), by considering \(k''=k\times k': Y \to P \times_{Z} Q\), we reduce to the case where \(Q\) dominates \(P\).
So we have a diagram such as the one shown, and we must show that the two \(c_{f,g}\) are compatible, i.e., that the following diagram is commutative:

\setcounter{page}{329}
\[f^{a} \stackrel{\gamma_{i,p}}{\to} i^{y} p^{z} \stackrel{IV}{\to} k^{y}(p r)^{z} \stackrel{IV}{\to} k^{y}(q s)^{z} \leftarrow II \to x_{s}^{z} q^{z} \stackrel{II}{\to} k^{y} q^{z} \stackrel{II}{\to} k^{y} p^{z} \stackrel{II}{\to} l^{y}(r t)^{z}=l^{y}(q s)^{z}<\Pi l^{y}(s t)^{z}<\Pi l^{y}(s t)^{z}<\Pi l^{y}(s t)^{z}<\Pi l^{y}(s t)^{z}<\Pi l^{y}(s t)^{z}<\Pi l^{y}(s t)^{z}<\Pi l^{y}(s t)^{z} q^{z}<\mathbb{N}} j^{y} q^{z}<\stackrel{\gamma}{\to} \ell^{y} .\]

One may verify this (~) using the definition of a permissible isomorphism, and the compatibilities ii, v, and vii above.

Similarly we must consider a change of the chart on \(X\) to another, say \(i': X \to \mathbb{A}_{Y}^{m}\) and as before we may assume that \(i'\) dominates \(i\).
This involves checking another analogous commutative diagram of isomorphisms, which we leave to the reader (!).

c) Construction of the isomorphism \(\psi_{f}\) for a finite morphism \(f\).
Let \(f: X \to Y\) be a finite morphism.
As before, it will be sufficient to construct the isomorphism locally on charts, provided our definition is compatible with permissible isomorphisms of charts.
So let \(i: X \to P\) be a chart for \(f\), and define

\setcounter{page}{330}
\[\psi_{f}: f^{a} \stackrel{\gamma_{i,p}}{\to} i^{y} p^{z} \stackrel{IV}{\leftarrow} f^{Y} .\]

If \(j: X \to Q\) is another chart, one may assume as usual that \(j\) dominates \(i\), and one checks (~) using compatibility \((v)\) above that \(\psi_{f}\) is compatible.

d) For a smooth morphism \(g: X \to Y\) we must construct the isomorphism \(\varphi_{g}\).
Take \(g\) as its own chart, and take \(\varphi_{g}=\gamma\)
There is no choice involved, hence nothing to check.

Having constructed the data a) - d), we must verify the conditions VAR 1 - VAR 5.
VAR 1). Given three morphisms f,g,h, we use a diagram such as the one shown to calculate \(c_{f,g}\), etc., on the charts.
The diagram of the condition becomes a large diagram whose commutative one checks (:) using (iii) and (iv) above.

VAR 2). If f and g are consecutive finite morphisms, one checks (~) that \(c_{f,g}\) is compatible with (I) via \(\psi_{f}\) and \(\psi_{g}\), by using (v) and (vi) above.

\setcounter{page}{331}
VAR 3). Trivial. One has only to observe that the isomorphism (III) is the identity if one side of the square is the identity.

VAR 4) and VAR 5) follow (:) from the definitions and a few more commutative diagrams.

This completes the proof of existence of the theory of variance, and we now show its uniqueness.
For that purpose, suppose that \(\{f^{\Delta}, c_{f,g}^{\Delta}, \psi_{f}^{\Delta}, \varphi_{g}^{\Delta}\}\) and \(\{f^{x}, c_{g}^{x}, \psi_{f}^{x}, \varphi_{g}^{x}\}\)
are two sets of data a)-d), each satisfying the conditions VAR 1 - VAR 5 (which we will call VAR \(1^{\Delta}\), VAR \(1^{x}\), etc.)
We will construct an isomorphism \(\delta: f^{\Delta} \to f^{x}\), compatible with the data b)-d) and b')-d'), and we will observe that \(\delta\) is unique.

Let \(f: X \to Y\) be a morphism of finite type, and choose locally an embedding \(i: X \to P\) into a scheme smooth over Y.
Define \(\delta\) by

\setcounter{page}{332}
\[\delta: f^{\Delta} \stackrel{c_{i,p}^{\Delta}}{\to} i^{\Delta} p^{\Delta} \stackrel{\psi_{i}^{\Delta}, \varphi_{p}^{\Delta}}{\to} i^{y} p^{z} \stackrel{\psi_{i}^{x}, \varphi_{p}^{x}}{\to} i^{x} p^{x} \stackrel{c_{i,p}^{x}}{\to} f^{x} .\]

(Note that in order to be compatible with \(c_{i,p}\), \(\psi_{i}\), \(\varphi_{p}\), we must choose \(\delta\) this way, which proves the uniqueness.)

To see that \(\delta\) is independent of the embedding chosen, let \(j: X \to Q\) be another one.
Replacing Q by \(P\times Q\) as usual we reduce to the case Q dominates P, and have a diagram such as the one shown.
We must check that the perimeter of the following diagram is commutative:

\setcounter{page}{333}
Thus \(\delta\) is well-defined locally. Now cover \(X\) by open sets \(U_{\nu}\) for which this is possible, and glue together the local isomorphisms
\[\delta_{\nu}:\left(f_{U_{\nu}}\right)^{\Delta} \stackrel{\sim}{\to}\left(f_{U_{\nu}}\right)^{x}\]
where \(j_{\nu}: U_{\nu} \to X\) is the open immersion.
Since \(j_{\nu}^{\Delta}\) is isomorphic via \(\varphi_{j_{\nu}}^{\Delta}\) with \(j_{\nu}^{z}\) which is the restriction, we can glue once we have checked (:) that the isomorphisms are compatible with restriction.
Thus
\[\delta: f^{\Delta} \to f^{x}\]
is defined.

Now we must check that \(\delta\) is compatible with the isomorphisms \(c_{f,g}\), \(\psi_{f}\), \(\varphi_{g}\).
This can be done locally using VAR 1 - VAR 5, and we leave the details to the reader (~). q.e.d.

\begin{proposition}
Let \(f: X \to Y\) be a morphism (with the conventions above) and let \(K^{\cdot}\) be a residual complex on \(Y\).
Let \(d\) denote the codimension function on \(X\) (resp. \(Y\)) associated with the pointwise dualizing complex \(f^{\Delta} K^{\cdot}\) (resp. \(K^{\cdot}\)).
Then for each \(x \in X\), if \(y=f(x)\), we have
\[d(x)=d(y)+\operatorname{tr.deg}k(x)/k(y) .\]
\end{proposition}

\setcounter{page}{334}
Proof. The question is local, and compatible with composition of morphisms.
Thus we reduce to the case \(f\) finite (trivial) or \(f\) smooth (which is [V 8.4]).

\begin{corollary}
Let \(f: X \to Y\) be a morphism (with the conventions above: \(f\) is of finite type, and the dimensions of its fibres are bounded).
Then if \(Y\) admits a residual complex (resp. dualizing complex) so does \(X\).
\end{corollary}

Proof. The existence of a residual complex on \(X\), given one on \(Y\), follows from the theorem.
If \(Y\) admits a dualizing complex \(R^{\cdot}\), then the associated codimension function \(d\) on \(Y\) is bounded.
By the proposition, the codimension function associated to \(Q f^{\Delta} E(R^{\cdot})\) is also bounded.
But \(Q f^{\Delta} E(R^{\cdot})\) is pointwise dualizing, hence dualizing (cf. proof of [V 8]).

\end{document}